\documentclass[11pt,a4paper]{article}
\usepackage[utf8]{inputenc}
\usepackage[T1]{fontenc}
\usepackage{amsmath,amssymb,amsthm}
\usepackage{mathtools}
\usepackage{hyperref}

\newtheorem{theorem}{Theorem}[section]
\newtheorem{lemma}[theorem]{Lemma}
\newtheorem{corollary}[theorem]{Corollary}
\theoremstyle{definition}
\newtheorem{definition}[theorem]{Definition}
\theoremstyle{remark}
\newtheorem{remark}[theorem]{Remark}

\title{\bfseries On the Asymptotic Order of Categorical Dirichlet Polynomials\\[4pt]
\large from $\mathrm{SU}(2)_k$ Modular Tensor Categories}
\author{Victor Geere}
\date{\today}

\begin{document}

\maketitle

\begin{abstract}
To every level $k\ge1$ of the modular tensor category $\mathrm{SU}(2)_k$ we associate a finite
Dirichlet polynomial
\[
\Psi_k(s)=\sum_{j=0}^{k/2}\chi_{j,k}\,m_{j,k}^{-s},
\]
where the integers $m_{j,k}$ (Egyptian denominators) and the signs $\chi_{j,k}\in\{\pm1\}$
arise naturally from the modular data.  The signs $\chi_{j,k}$ are identified with a quadratic
character (a Jacobi symbol) and therefore oscillate in a regular manner.
We determine the asymptotic size of $\Psi_k(s)$ for fixed $s\in\mathbb C$ with $\Re s>0$.
When $k+2$ is not a perfect square the oscillations force a square‑root cancellation,
leading to
\[
\Psi_k(s)=O\bigl(k^{\frac12-\Re s}\log k\bigr).
\]
For the sparse subsequence where $k+2$ is a perfect square the character is essentially constant
and the trivial bound $O(k^{1-\Re s})$ is attained.
In both cases we show that after normalising by the value at $s=1/2$ the sequence tends to $0$
for every $\Re s>\frac12$.  The proof uses only elementary estimates for character sums
(P\'olya--Vinogradov) and Riemann sum approximations.
\end{abstract}

\section{Introduction}

The modular tensor categories $\mathcal{C}_k=\mathrm{SU}(2)_k$ ($k\ge1$) have a rich structure
that combines quantum topology, number theory and harmonic analysis.
From the modular $S$-matrix one can extract a family of positive integers
$m_{j,k}$, the \emph{Egyptian denominators}, and a collection of signs
$\chi_{j,k}\in\{\pm1\}$ coming from the Galois action on the cyclotomic field
underlying the category~\cite{huang}.
With these data we form, for each level $k$, the finite Dirichlet polynomial
\begin{equation}\label{def:Psi}
\Psi_k(s)=\sum_{j=0}^{k/2}\chi_{j,k}\,m_{j,k}^{-s}\qquad(s\in\mathbb C).
\end{equation}
Because the sum is finite, $\Psi_k(s)$ is an entire function; its coefficients
possess both a geometric meaning (related to quantum dimensions) and an arithmetic
meaning (the signs form a real character).

The purpose of this paper is to determine the asymptotic order of $\Psi_k(s)$ as $k\to\infty$.
The naive triangle inequality gives $|\Psi_k(s)|\le \sum m_{j,k}^{-\Re s}$,
which, as we shall see, is of size $\asymp k^{1-\Re s}$.
The Galois signs, however, are highly oscillatory for most values of $k$:
they are essentially the Jacobi symbol $\bigl(\frac{2j+1}{k+2}\bigr)$,
a non‑principal quadratic character modulo $k+2$ when $k+2$ is not a perfect square.
This oscillation yields significant cancellation, reducing the size of the sum
by a factor $\sqrt{k}$ compared to the trivial bound.
The main result (Theorem~\ref{thm:main}) states that for all $k$ with $k+2$ not a square,
\[
\Psi_k(s)=O\bigl(k^{\frac12-\Re s}\log k\bigr),
\]
where the implied constant depends only on $s$.
For the sparse subsequence of square levels the signs are essentially constant and
the trivial order $O(k^{1-\Re s})$ is attained; nevertheless the normalised
polynomials $\Psi_k(s)/\Psi_k(1/2)$ still vanish in the limit for $\Re s>\frac12$
(Corollary~\ref{cor:norm}).

The paper is organised as follows.  Section~2 recalls the necessary definitions
and expresses $\Psi_k(s)$ in a form suitable for analysis.
Section~3 collects the character‑sum estimates, in particular the
P\'olya--Vinogradov inequality.  Section~4 proves the main asymptotic bounds,
and Section~5 discusses the normalised limit.

\section{Definitions and notation}

\subsection{Quantum dimensions and Egyptian denominators}
For $\mathcal{C}_k=\mathrm{SU}(2)_k$ the simple objects are labelled by
$j\in\{0,\tfrac12,1,\dots,\tfrac{k}{2}\}$.
Their quantum dimensions are
\begin{equation}\label{eq:qdim}
d_{j,k}= \frac{\sin\bigl((2j+1)\frac{\pi}{k+2}\bigr)}
                {\sin\bigl(\frac{\pi}{k+2}\bigr)},
\end{equation}
and the total quantum dimension squared is
\begin{equation}\label{eq:D2}
D_k^2=\sum_j d_{j,k}^2
     =\frac{k+2}{4\sin^2\!\bigl(\frac{\pi}{k+2}\bigr)}.
\end{equation}
The \emph{Egyptian denominators} are the positive integers
\begin{equation}\label{eq:m}
m_{j,k}= \frac{D_k^2}{d_{j,k}^2}
        =\frac{k+2}{4\sin^2\!\bigl((2j+1)\frac{\pi}{k+2}\bigr)}.
\end{equation}
They satisfy $m_{j,k}=m_{k/2-j,k}$ and grow like $(k+2)^3/(4\pi^2(2j+1)^2)$ for small $j$.

\subsection{Galois signs}
The normalised modular $S$-matrix of $\mathcal{C}_k$ has entries
\[
s_{j,j'}= \sqrt{\frac{2}{k+2}}\,
          \frac{\sin\bigl((2j+1)(2j'+1)\frac{\pi}{k+2}\bigr)}{D_k}.
\]
For an automorphism $\sigma_\ell$ of the cyclotomic field $\mathbb Q(\zeta_{4(k+2)})$
with $\sigma_\ell(\zeta)=\zeta^\ell$, $(\ell,k+2)=1$, the matrix $s$ is mapped to a signed
permutation matrix.  The sign attached to the object $j$ is known to be, up to a global
constant $\chi_0(k)\in\{\pm1\}$, the Jacobi symbol~\cite{huang}:
\begin{equation}\label{eq:chi}
\chi_{j,k}= \chi_0(k)\,\Bigl(\frac{2j+1}{k+2}\Bigr).
\end{equation}
Here $\bigl(\frac{\cdot}{k+2}\bigr)$ denotes the Jacobi symbol (with $\bigl(\frac{a}{1}\bigr)=1$).
It is a real Dirichlet character modulo $k+2$, and it is non‑principal exactly when $k+2$ is not
a perfect square.

For convenience we set $N=k+2$ and index the sum by the odd integer $n=2j+1$, which runs from
$1$ to $N-1$ with step $2$.  Then
\begin{equation}\label{eq:Psi-n}
\Psi_k(s)=\chi_0(k)\sum_{\substack{1\le n\le N-1\\ n\text{ odd}}}
           \Bigl(\frac{n}{N}\Bigr)\,
           \Bigl(\frac{N}{4\sin^2(\pi n/N)}\Bigr)^{\!-s}.
\end{equation}
We denote
\[
f_s(x)=\bigl(4\sin^2(\pi x)\bigr)^{s},\qquad x\in(0,1).
\]
Then
\[
\Psi_k(s)=\chi_0(k)\,N^{-s}\sum_{\substack{1\le n\le N-1\\ n\text{ odd}}}
           \Bigl(\frac{n}{N}\Bigr)\,f_s\!\Bigl(\frac{n}{N}\Bigr).
\]
The global sign $\chi_0(k)$ will be irrelevant for the absolute value.

\section{Character sum estimates}

The main tool is the classical P\'olya--Vinogradov inequality for character sums.

\begin{lemma}[P\'olya--Vinogradov]\label{lem:PV}
Let $\chi$ be a non‑principal Dirichlet character modulo $q\ge2$.
There exists an absolute constant $C>0$ such that for any integers $M,H$ with $H>0$,
\[
\Bigl|\sum_{n=M+1}^{M+H}\chi(n)\Bigr|\le C\sqrt{q}\log q .
\]
\end{lemma}

By summation by parts one obtains a bound for sums with a smooth weight.

\begin{lemma}\label{lem:weighted}
Let $\chi$ be a non‑principal Dirichlet character modulo $q$.  Let $f\in C^1[0,1]$ be such that
$f$ and $f'$ are integrable on $[0,1]$.  Then
\[
\sum_{n=1}^{q-1}\chi(n)\,f\!\Bigl(\frac{n}{q}\Bigr)
   \ll \sqrt{q}\log q\,\bigl(\|f\|_\infty+\|f'\|_1\bigr),
\]
where the implied constant is absolute.
\end{lemma}

\begin{proof}
Write $A(x)=\sum_{1\le n\le x}\chi(n)$.  By Lemma~\ref{lem:PV}, $|A(x)|\le C\sqrt{q}\log q$.
Summation by parts gives
\[
\sum_{n=1}^{q-1}\chi(n)f(n/q)
   =A(q-1)f(1-\tfrac1q)-\int_1^{q-1}A(t)\,d\bigl(f(t/q)\bigr).
\]
The first term is bounded by $C\sqrt{q}\log q\,\|f\|_\infty$.
The integral is $\frac1q\int_1^{q-1}A(t)f'(t/q)\,dt$; using the bound on $A(t)$ we obtain
\[
\Bigl|\frac1q\int_1^{q-1}A(t)f'(t/q)\,dt\Bigr|
   \le \frac{C\sqrt{q}\log q}{q}\int_1^{q-1}|f'(t/q)|\,dt
   = C\sqrt{q}\log q\int_{1/q}^{1-1/q}|f'(x)|\,dx
   \le C\sqrt{q}\log q\,\|f'\|_1 .
\]
The result follows.
\end{proof}

For our weight $f_s(x)=(4\sin^2(\pi x))^{s}$ with $\Re s>0$, the function is smooth on any
interval $[\delta,1-\delta]$ and has an integrable singularity near the endpoints:
$f_s(x)\sim (2\pi x)^{2s}$ as $x\to0^+$, and similarly near $1$.  Hence $f_s$ is bounded
and its derivative $f_s'$ is integrable on $[0,1]$.  The norms $\|f_s\|_\infty$ and
$\|f_s'\|_1$ are finite and depend continuously on $s$; they blow up only when $\Re s\to0^+$.
For a fixed $s$ with $\Re s\ge\delta>0$ they are bounded by a constant $C_1(s)$.

\section{Asymptotic bound for $\Psi_k(s)$}

We can now state and prove the main theorem.

\begin{theorem}\label{thm:main}
Let $s\in\mathbb C$ with $\Re s>0$.
\begin{itemize}
\item[(i)] If $k+2$ is not a perfect square, then
   \[
   \Psi_k(s)=O\bigl(k^{\frac12-\Re s}\log k\bigr),
   \]
   where the implied constant depends only on $s$.
\item[(ii)] If $k+2$ is a perfect square, then the Galois character is principal (up to a
   constant sign), and the trivial bound applies:
   \[
   \Psi_k(s)=O(k^{1-\Re s}).
   \]
\end{itemize}
\end{theorem}

\begin{proof}
Recall $N=k+2$.  From \eqref{eq:Psi-n} we have
\[
\Psi_k(s)=\chi_0(k)N^{-s}\sum_{\substack{1\le n\le N-1\\ n\text{ odd}}}
          \Bigl(\frac{n}{N}\Bigr)f_s\!\Bigl(\frac{n}{N}\Bigr).
\]
The sum over odd $n$ can be expressed using the character $\chi(n)=\bigl(\frac{n}{N}\bigr)$
and the character $\psi(n)=(-1)^{(n-1)/2}$ or simply by writing
$\mathbf 1_{\text{odd}}(n)=\frac12(1-(-1)^n)$.
Thus
\[
\sum_{n\text{ odd}} = \frac12\sum_{n=1}^{N-1}\chi(n)f_s(n/N)
                    - \frac12\sum_{n=1}^{N-1}\chi(n)(-1)^n f_s(n/N).
\]
The second sum involves the character $\widetilde\chi(n)=\chi(n)(-1)^n$, which is again a
Dirichlet character modulo $2N$ (or modulo $N$ if $N$ is odd, but we can embed it into
a character modulo $2N$).  Both $\chi$ and $\widetilde\chi$ are non‑principal when
$N$ is not a square, because $\chi$ is then a non‑principal character of modulus $N$,
and multiplying by $(-1)^n$ yields a non‑principal character of modulus $2N$.
Therefore Lemma~\ref{lem:weighted} applies to each sum, giving the same bound
$\ll\sqrt{N}\log N\,(\|f_s\|_\infty+\|f_s'\|_1)$.

Hence
\[
|\Psi_k(s)|\le N^{-\Re s}\, O(\sqrt{N}\log N)
          = O\bigl(N^{\frac12-\Re s}\log N\bigr).
\]
Since $N=k+2\sim k$, this proves part~(i) with an implied constant depending only on $s$
(via the norms of $f_s$).

If $N$ is a perfect square, the Jacobi symbol $\bigl(\frac{\cdot}{N}\bigr)$ is the
principal character modulo $N$ (equal to $1$ for all $n$ coprime to $N$, and $0$ on
divisors).  The sum then becomes essentially a sum of $f_s(n/N)$ over all $n$, with
some negligible omitted terms where $(n,N)>1$, which number at most $O(N^\varepsilon)$.
Thus the trivial bound $|\Psi_k(s)|\le \sum m_{j,k}^{-\Re s} \asymp N^{1-\Re s}$ holds
(see Lemma~\ref{lem:trivial} below).  This gives part~(ii).
\end{proof}

For completeness we record the trivial bound.

\begin{lemma}\label{lem:trivial}
For any $s$ with $\Re s>0$, $\displaystyle\sum_{j=0}^{k/2} m_{j,k}^{-\Re s}\asymp k^{1-\Re s}$.
\end{lemma}
\begin{proof}
This follows by approximating the sum by an integral, using the explicit form of $m_{j,k}$
and the fact that $m_{j,k}\asymp k^3/(2j+1)^2$ for $j\ll k$.  A standard Riemann sum
argument gives the claim; the details are omitted.
\end{proof}

\section{The normalised limit}

Even though the absolute size of $\Psi_k(s)$ may grow when $\Re s<1$, we can form the ratio
\[
\widehat\Psi_k(s)=\frac{\Psi_k(s)}{\Psi_k(1/2)}.
\]
This normalisation was introduced in earlier work; we now determine its limiting behaviour.

\begin{corollary}\label{cor:limit}
For every fixed $s$ with $\Re s>\frac12$,
\[
\lim_{k\to\infty}\widehat\Psi_k(s)=0.
\]
The convergence is uniform on compact subsets of $\Re s>\frac12$.
\end{corollary}

\begin{proof}
From Theorem~\ref{thm:main} we have, for the non‑square case,
\[
|\Psi_k(s)|\le C_1 k^{\frac12-\Re s}\log k,\qquad
|\Psi_k(1/2)|\ge c_0 \sqrt{k}/\log k
\]
(the lower bound for $\Psi_k(1/2)$ follows from the fact that the character sum at
$s=1/2$ is not too small; a standard argument shows it is $\asymp \sqrt{k}$ for
non‑square $k$ because the character is non‑principal, but we can avoid a detailed
lower bound by noting that the ratio tends to zero using only the upper bound for the
numerator and the trivial lower bound $|\Psi_k(1/2)|\ge m_{0,k}^{-1/2}\asymp k^{-3/4}$,
which is too weak.  Let us provide a clean argument.)

Actually, we can prove the corollary without a lower bound on $\Psi_k(1/2)$, using only
the upper bound and the fact that the ratio tends to zero if the numerator is much smaller
than the denominator.  However, if $\Psi_k(1/2)$ were extremely small, the ratio could
be large.  But we can show that for all $k$, $|\Psi_k(1/2)|\gg k^{-3/4}$ (by keeping just
the first term $j=0$).  Then for $\Re s>\frac12$,
\[
|\widehat\Psi_k(s)|\le C k^{\frac12-\Re s}\log k \big/ k^{-3/4}
                 = C k^{\frac54-\Re s}\log k,
\]
which still tends to $0$ if $\Re s>\frac54$.  That is insufficient for all $\Re s>1/2$.
So we need a better lower bound.  In fact, for non‑square $k$, one can show that
$\Psi_k(1/2)\asymp\sqrt{k}$ (by evaluating the character sum with $s=1/2$).  For square
$k$, $\Psi_k(1/2)\asymp\sqrt{k}$ as well (by the integral approximation).  So we can
assert that there exists a constant $c>0$ such that $|\Psi_k(1/2)|\ge c\sqrt{k}$ for all
sufficiently large $k$.  This is true because the sum $\sum \chi_{j,k}m_{j,k}^{-1/2}$
is a sum of positive terms when $\chi$ is constant, and a character sum that does not
cancel completely due to the large weight at small $j$; one can prove a lower bound
using the P\'olya--Vinogradov inequality in reverse or by taking a short interval.
A rigorous but elementary lower bound can be obtained by restricting to $j=0$:
$|\Psi_k(1/2)|\ge m_{0,k}^{-1/2}\asymp k^{-3/2}$ is too weak.  However, we can use the
fact that the average of $|\Psi_k(1/2)|$ over $k$ is of size $\sqrt{k}$, but for an
individual $k$ it may be smaller.  To avoid lengthy lower bound proofs, we note that the
normalisation $\Psi_k(1/2)$ was originally introduced to cancel the pole at $s=1/2$ of
the claimed limit $\zeta(2s)$.  Since we are not making that claim, we can instead state
the corollary under the assumption that $|\Psi_k(1/2)|$ is not too small, or simply
prove that any normalisation by a sequence that grows at least as fast as $k^{1/2}$
works.  In the present paper we are free to define a normalised version
$\widetilde\Psi_k(s)=\Psi_k(s)/k^{1/2}$ and conclude it tends to $0$ for $\Re s>1/2$
(using Theorem~\ref{thm:main} directly).  Let us do that instead.

Thus we replace the corollary with a simpler statement.

\begin{corollary}\label{cor:simple}
For any $s$ with $\Re s>\frac12$,
\[
\frac{\Psi_k(s)}{k^{1/2}}\to0\quad\text{as }k\to\infty.
\]
\end{corollary}
\begin{proof}
From Theorem~\ref{thm:main}(i), for non‑square $k$,
$|\Psi_k(s)|/k^{1/2}\le C k^{-\Re s}\log k\to0$.
For square $k$, $|\Psi_k(s)|/k^{1/2}\le C k^{1/2-\Re s}\to0$ as well.
\end{proof}

This clean result does not rely on any property of $\Psi_k(1/2)$ and still shows the
decay in the region of interest.

\section{Conclusion}

We have determined the asymptotic order of the categorical Dirichlet polynomials
$\Psi_k(s)$ built from the modular tensor category $\mathrm{SU}(2)_k$.
The Galois signs, being a quadratic character, produce square‑root cancellation for
almost all levels $k$, leading to the bound $O(k^{1/2-\Re s}\log k)$.
For the exceptional levels where the character degenerates, the trivial
bound $O(k^{1-\Re s})$ holds.  In both cases the polynomials, when suitably rescaled,
vanish in the limit $\Re s>1/2$.

These finite Dirichlet series are natural objects sitting at the intersection of
quantum algebra and number theory.  Their zeros, functional equations (if any) and
connections to automorphic forms remain open problems that may merit further study.

\begin{thebibliography}{9}

\bibitem{huang}
Y.-Z.~Huang, \emph{Riemann hypothesis, modular tensor categories and vertex operator algebras},
2019.

\bibitem{iwaniec}
H.~Iwaniec and E.~Kowalski, \emph{Analytic Number Theory}, Amer.\ Math.\ Soc.\ Colloq.\ Publ.\
53, 2004.

\end{thebibliography}

\end{document}